\section{Related Works}

\subsection{Person Representation Learning}
Early approaches to text-based person retrieval typically employ separate vision and language encoders with custom alignment losses~\cite{zheng2020dual, si2018dual}. These methods often exhibit suboptimal modality alignment and require extensive manual annotation. The introduction of CLIP~\cite{radford2021learning} establishes a unified vision-language embedding space, significantly advancing cross-modal matching. Recent works extend CLIP with specialized modules for text-based person retrieval. IRRA~\cite{jiang2023cross} merges visual cues into textual tokens via a cross-modal transformer and aligns global similarity distributions. MDRL~\cite{yang2024multi} designs a cross-modality global feature learning architecture to learn the global features from the two modalities and meet the demand of the task. UniPT~\cite{shao2023unified} utilizes a simple vision-and-language pre-training framework to explicitly align the feature space of the image and text modalities during pre-training. However, these approaches largely ignore data noise, which critically influences cross-modal alignment in feature space. RDE~\cite{qin2024noisy} mitigates the adverse impact of noisy through the proposed confident consensus division and novel triplet alignment loss. ProPOT~\cite{yan2024prototypical} transforms the identity-level matching problem into a prototype learning problem, aiming to learn identity-enriched prototypes. However, prototype aggregation compromises fine-grained semantic learning.


\subsection{Person-centric Dataset}
High-quality image-text paired datasets are essential for learning discriminative person representations. However, existing manually annotated datasets (e.g., CUHK-PEDES \cite{li2017person}, ICFG-PEDES \cite{ding2021semantically}, RSTPReid \cite{zhu2021dssl}) face severe scalability limitations due to labor-intensive annotation processes. This scalability bottleneck ultimately constrains models' capacity to acquire diverse semantic information and learn discriminative features. Recent efforts to mitigate this issue focus on constructing large-scale datasets, such as LUPerson \cite{fu2021unsupervised}, LUPerson-T \cite{shao2023unified}, LUPerson-MLLM \cite{tan2024harnessing}, and SYNTH-PEDES \cite{zuo2024plip} demonstrate that increased data volume improves general pedestrian representation learning. Nevertheless, these datasets primarily derive from video sources, inheriting inherent scalability constraints from computationally intensive video processing pipelines. The success of multimodal large language models in cross-modal understanding \cite{yu2024capsfusion} has inspired their application to synthetic data generation. For instance, LUPerson-MLLM \cite{tan2024harnessing} employs template-guided MLLMs to generate diverse textual descriptions, significantly enhancing text-to-image ReID performance. However, this approach remains limited by its dependence on existing LUPerson image collections.

